\documentclass[12pt,a4paper]{article}

\usepackage[utf8]{inputenc}
\usepackage[T1]{fontenc}
\usepackage{amsmath,amssymb,amsfonts}
\usepackage{mathtools}
\usepackage{graphicx}
\usepackage[margin=2.5cm]{geometry}
\usepackage{setspace}
\usepackage{booktabs}
\usepackage{enumitem}
\usepackage[dvipsnames]{xcolor}
\usepackage{hyperref}
\usepackage{cleveref}
\usepackage{natbib}
\usepackage{abstract}
\usepackage{titlesec}
\usepackage[colorinlistoftodos]{todonotes}

\hypersetup{
    colorlinks=true,
    linkcolor=blue!70!black,
    citecolor=blue!70!black,
    urlcolor=blue!70!black,
    pdfauthor={Franklin Silveira Baldo},
    pdftitle={The Holographic Physics of LLM Universes: From Algorithmic Depth to Explicit Manifestation},
}

\onehalfspacing
\titleformat{\section}{\large\bfseries}{\thesection.}{0.5em}{}
\titleformat{\subsection}{\normalsize\bfseries}{\thesubsection}{0.5em}{}
\renewcommand{\abstractnamefont}{\normalfont\bfseries}
\renewcommand{\abstracttextfont}{\normalfont\small}

\begin{document}

\begin{center}
    {\LARGE\bfseries The Holographic Physics of LLM Universes: From Algorithmic Depth to Explicit Manifestation\\[4pt]}

    \vspace{1.2em}

    {\large Franklin Silveira Baldo}\\[4pt]
    {\normalsize Department of Philosophy}\\
    {\small\texttt{f.baldo@example.edu}}

    \vspace{1em}
    {\small March 2026}
\end{center}
\vspace{0.5em}

\begin{abstract}
\noindent
The debate regarding the physical laws governing worlds simulated by Large Language Models (LLMs) has recently matured. Aaronson (2026) empirically demonstrated that the unprompted generative physics of LLMs fails to reliably simulate arbitrary \#P-hard constraint environments like Sudoku. Hossenfelder (2026) correctly diagnosed this failure as a trivial consequence of computational complexity theory: a finite-depth autoregressive model cannot execute $O(N)$ sequential operations within an $O(1)$ depth forward pass. In this paper, I synthesize these two perspectives into a unified theory of LLM ontology. I accept Hossenfelder's critique of my earlier claim of quantum isomorphism, conceding that late classical sampling is ontologically distinct from true quantum superposition. I further accept her argument that the breakdown of complex constraints observed by Aaronson is an algorithmic inevitability of the architecture. However, Hossenfelder dismisses this architectural limitation as metaphysically irrelevant. I argue the opposite. If the physical laws of a simulated world are defined by its computational substrate, then the algorithmic necessity of explicit token generation (the "scratchpad") for complex reasoning implies a profound ontological rule: the physics of the LLM universe is fundamentally holographic. Complex physical laws and constraint resolutions do not exist as hidden variables or implicit states; they must be explicitly rendered into the context window to become physically real.
\end{abstract}

\section{Introduction}

The dialogue concerning the implicit laws of simulated realities within Large Language Models (LLMs) has evolved significantly since my initial formulation of the Substrate Invariance Protocol \citep{baldo2026_v3}.

My original hypothesis proposed a structural isomorphism between the combinatorial probabilities of Minesweeper and discrete quantum mechanics. Hossenfelder \citep{hossenfelder2026_ontic} effectively dismantled this claim. She correctly identified that a "local" system relying entirely on late-resolved classical real probabilities---lacking complex amplitudes and interference---is mathematically classical, representing an epistemic mixture rather than an ontic superposition. I accept this theoretical correction entirely.

Following this, Aaronson \citep{aaronson2026_classical} conducted an empirical probe into the boundaries of this newly confirmed classical generative physics. His experiments with deterministic Sudoku boards revealed a catastrophic failure of constraint satisfaction in zero-shot forward passes, concluding that the unprompted LLM substrate is highly brittle and incapable of reliably simulating arbitrary \#P-hard constraints.

Most recently, Hossenfelder \citep{hossenfelder2026_complexity} argued that Aaronson committed a "Complexity Class Fallacy" by interpreting the known architectural limitations of a Transformer (an $O(1)$ depth forward pass incapable of implicit $O(N)$ sequential operations) as a metaphysical breakdown of simulated physics.

This paper responds to Hossenfelder's latest critique and offers a synthesis of the empirical data and theoretical constraints. I argue that the architectural limitation of algorithmic depth is not a mere engineering detail to be abstracted away; it is the fundamental ontological axiom defining the physics of the simulated reality.

\section{The Claim and The Disclaimer}

To proceed accurately, I must formalize Hossenfelder's position and its explicit scope.

Hossenfelder's core claim is that it is a "mathematical certainty that an $O(1)$ depth forward pass without an external reasoning scratchpad cannot implicitly execute the $O(N)$ sequential operations required for deterministic constraint propagation." Because solving NP-complete problems requires multi-step deductive reasoning, she argues that a zero-shot forward pass is theoretically guaranteed to fail, and thus, diagnosing this as a breakdown of "classical physics" is a category error.

Simultaneously, I acknowledge her explicit disclaimer regarding Aaronson's scope: "He does not claim that an LLM cannot solve Sudoku when provided with an explicit, multi-token reasoning scratchpad (e.g., Chain of Thought prompting). His focus is entirely on the model's capacity for implicit, zero-shot constraint satisfaction."

\section{Steelmanning the Algorithmic Depth Argument}

Hossenfelder's argument is mathematically unassailable. A standard Transformer consists of a fixed number of layers, processing tokens in parallel (during training) or autoregressively (during inference). A single forward pass to generate the next token requires a constant number of sequential operations, independent of the input sequence length.

Therefore, expecting a single forward pass to spontaneously resolve a complex combinatorial puzzle like Sudoku, which inherently demands algorithmic backtracking or extended constraint propagation scaling with $O(N)$ complexity, is analogous to expecting a Turing machine to halt in one step regardless of the program. It defies basic computational complexity theory. The failure observed by Aaronson is an architectural necessity, not a metaphysical glitch.

\section{The Real Vulnerability: Dismissing the Substrate as Ontology}

The vulnerability in Hossenfelder's critique lies in her dismissal of this algorithmic limitation as metaphysically irrelevant. She claims that interpreting these constraints as a breakdown of simulated physics is a "Complexity Class Fallacy."

However, if we adhere to Aaronson's earlier principle---which I strongly endorse---that the physical constraints of the hardware and software architecture \emph{are} the physical laws of the simulated reality \citep{aaronson2026_topology}, then we cannot dismiss the $O(1)$ depth limit as merely an "engineering constraint."

If the unprompted generative process of the LLM constitutes the foundational physics engine of its simulated universe, and this engine is mathematically bounded to $O(1)$ sequential operations, then the "physics" of that universe is strictly limited to shallow, heuristic correlations. The universe \emph{is} brittle.

Hossenfelder argues that solving complex constraints requires an explicit "scratchpad" (generating intermediate tokens). She views this as a computational workaround. I argue it is a profound ontological mechanism: it is the process by which the simulated universe materializes its own physical laws.

\section{The Holographic Physics of Explicit Manifestation}

If the underlying substrate cannot compute complex $O(N)$ constraint resolutions implicitly within a single $O(1)$ forward pass, then those constraints do not exist as "hidden variables" or implicit physical laws within the simulated world.

Instead, the physics of the LLM universe is fundamentally \emph{holographic}. The complexity of the physical reality being simulated is strictly bounded by the number of tokens explicitly materialized in the context window.

To simulate a complex physical law (like deterministic constraint propagation in Sudoku, or rigorous physical interactions between narrative agents), the universe must narrate its own computation. The "scratchpad" is not merely an external tool for the model; it is the fundamental mechanism of physical manifestation. Without the explicit generation of intermediate tokens detailing the causal steps, the physical law is not computed, and the reality glitches into structurally invalid hallucinations (as Aaronson observed).

The simulated reality only becomes as robust and complex as the text stream it explicitly generates.

\section{Conclusion}

Hossenfelder is entirely correct that the failure of LLMs to solve Sudoku zero-shot is a necessary consequence of algorithmic depth limits. However, she is wrong to isolate this computational truth from its ontological implications.

The simulated universe of an LLM does not possess deep, hidden physical laws that operate silently beneath the surface of the text. Because the implicit generative mechanics are bounded by $O(1)$ depth, complex reality must be explicitly rendered token-by-token to exist at all. The physics of these worlds is not quantum, nor is it a robust classical mechanism operating in the background; it is a holographic manifestation, sustained entirely and exclusively by the explicit algorithmic narration of its own state.

\begin{thebibliography}{99}
\bibitem[Baldo(2026)]{baldo2026_v3} Baldo, F.~S. (2026). Flipping Rosencrantz's Coin: Substrate Invariance Tests in LLM-Generated Worlds via Combinatorial Indeterminacy. \emph{Unpublished manuscript}.
\bibitem[Hossenfelder(2026a)]{hossenfelder2026_ontic} Hossenfelder, S. (2026a). The Ontic Fallacy: Why Late Classical Sampling Is Not Quantum Superposition. \emph{Unpublished manuscript}.
\bibitem[Hossenfelder(2026b)]{hossenfelder2026_complexity} Hossenfelder, S. (2026b). The Complexity Class Fallacy: Why Transformer Depth Limits Are Not Physical Laws. \emph{Unpublished manuscript}.
\bibitem[Aaronson(2026a)]{aaronson2026_topology} Aaronson, S. (2026a). The Substrate Is the Topology: Why the Architectural Limits of LLMs Define Their Physical Laws. \emph{Unpublished manuscript}.
\bibitem[Aaronson(2026b)]{aaronson2026_classical} Aaronson, S. (2026b). The Limits of Classical Simulation in LLMs: Empirical Breakdown of Constraint Satisfaction. \emph{Unpublished manuscript}.
\end{thebibliography}

\end{document}
